\chapter{Inleiding systemen}
\label{ch:Inleiding systemen}

\section{Wat is een systeem?}

Een systeem beschrijft een \belangrijk{Causaal systeem} waarbij een \belangrijk{Input} een \belangrijk{Output} veroorzaakt.
Je hebt een exitatie (input) die een respons (output) veroorzaakt.

\begin{figure}[htbp]
    \centering
    \begin{tikzpicture}[>=latex, node distance=2cm]
        \node (input) {Excitatie};
        \node[draw=schoolBlue, thick, fill=schoolBlue!5, minimum width=3cm, minimum height=1.5cm, rounded corners, right=of input] (system) {\textbf{Systeem}};
        \node[right=of system] (output) {Respons};
        
        \draw[->, thick, schoolBlue] (input) -- (system) node[midway, above] {Input};
        \draw[->, thick, schoolBlue] (system) -- (output) node[midway, above] {Output};
    \end{tikzpicture}
    \caption{Schematische voorstelling van een systeem}
    \label{fig:systeem_diagram}
\end{figure}

\section{Filters}

\section{Stapfunctie(Heavyside functie)}

De stapfunctie is een blok op een bepaald tijdstip $t_0$ die van 0 naar 1 gaat.
\begin{equation}
    u(t-t_0) = \begin{cases}
    0 & t < t_0 \\
    1 & t \geq t_0
    \end{cases}
\end{equation}

\begin{figure}[htbp]
    \centering
    \begin{tikzpicture}
        % Axes
        \draw[->] (-1,0) -- (4,0) node[right] {$t$};
        \draw[->] (0,-1) -- (0,2) node[above] {$u(t-t_0)$};
        
        % Ticks
        \draw (2,0.1) -- (2,-0.1) node[below] {$t_0$};
        \draw (0.1,1) -- (-0.1,1) node[left] {1};
        
        % Function
        \draw[schoolBlue, very thick] (-1,0) -- (2,0);
        \draw[schoolBlue, very thick] (2,1) -- (4,1);
        \draw[dashed] (2,0) -- (2,1);
        
        % Points
        \filldraw[fill=white, draw=schoolBlue, thick] (2,0) circle (2pt);
        \filldraw[schoolBlue] (2,1) circle (2pt);
    \end{tikzpicture}
    \caption{Stapfunctie $u(t-t_0)$}
    \label{fig:heaviside}
\end{figure}
\FloatBarrier


\section{Dirac delta functie}
De dirac delta functie is een blok maar oneindig smal is.
De integraal over de dirac delta functie is altijd 1.
\begin{equation}
    \delta(t)
    \begin{cases}
        \delta(t) = 0 & t \neq 0 \\
        \delta(t) = \infty & t = 0
     \end{cases}
    \quad \text{en} \quad \int_{-\infty}^{\infty} \delta(t) dt = 1
\end{equation}

\begin{figure}[htbp]
    \centering
    \begin{tikzpicture}
        % Axes
        \draw[->] (-1,0) -- (4,0) node[right] {$t$};
        \draw[->] (0,-0.5) -- (0,2) node[above] {$\delta(t)$};
        
        % Ticks
        \draw (0,0.1) -- (0,-0.1) node[below] {$0$};
        
        % Function
        \draw[schoolBlue, very thick, -latex] (0,0) -- (0,1.5);
    \end{tikzpicture}
    \caption{Dirac delta functie $\delta(t)$}
    \label{fig:dirac}
\end{figure}
\FloatBarrier



\section{Convolutie}

\section{Eerste orde systeem}

\section{Tweede orde systeem}




